% 子集（综述）
% license CCBYSA3
% type Wiki

本文根据 CC-BY-SA 协议转载翻译自维基百科\href{https://en.wikipedia.org/wiki/Subset}{相关文章}

\begin{figure}[ht]
\centering
\includegraphics[width=6cm]{./figures/d095c5b69171d4da.png}
\caption{欧拉图示意：A 是 B 的子集（记作 $A \subseteq B$），反之，B 是 A 的超集（记作 $B \supseteq A$）。} \label{fig_ZiJ_1}
\end{figure}
在数学中，如果集合A的所有元素也是集合B的元素，则称A是B的子集；此时$B$是$A$的超集。集合$A$与$B$有可能相等；如果它们不相等，则A是B的真子集。一个集合是另一个集合的子集这一关系称为包含关系（有时也称为包涵关系）。$A$是$B$的子集也可以表述为$B$包含$A$或$A$被包含于$B$。一个k-子集是指具有k个元素的子集。

在量化表示时，$A \subseteq B$可写作$\forall x \,(x \in A \Rightarrow x \in B)$。

可以通过一种称为元素论证法的证明技巧来证明命题$A \subseteq B$。
\begin{enumerate}
\item 设给定集合$A$和$B$。要证明$A \subseteq B$。
\item 假设a是集合A中一个特定但任意选择的元素，并证明a也是集合B的元素。
\end{enumerate}
这种方法的有效性可以看作是全称概化的结果：该方法表明，对于任意选择的元素c，都有
$(c \in A) \Rightarrow (c \in B)$。由此，全称概化推出$\forall x \,(x \in A \Rightarrow x \in B)$，这等价于$A \subseteq B$如上所述。
\subsection{定义}
如果$A$和$B$是集合，且$A$的每一个元素也是$B$的元素，那么：
\begin{itemize}
\item $A$是$B$的子集，记作$A \subseteq B$或等价地，
\item $B$是$A$的超集，记作$B \supseteq A$
\end{itemize}
如果A是B的子集，但A不等于B（即至少存在一个B中的元素不在A中），那么：
\begin{itemize}
\item $A$是$B$的真（或严格）子集，记作$A \subsetneq B$或等价地，
\item $B$是$A$的真（或严格）超集，记作$B \supsetneq A$。
\end{itemize}
空集，写作$\{\}$或$\varnothing$，没有任何元素，因此平凡地是任意集合$X$的子集。
\subsection{基本性质}
\begin{figure}[ht]
\centering
\includegraphics[width=6cm]{./figures/241b397d38ab5c05.png}
\caption{$A \subseteq B \quad \text{且} \quad B \subseteq C \;\;\Rightarrow\;\; A \subseteq C$} \label{fig_ZiJ_2}
\end{figure}
\begin{itemize}
\item 自反性：对任意集合$A$都有$A \subseteq A$
\item 传递性：若$A \subseteq B$且$B \subseteq C$则$A \subseteq C$
\item 反对称性：若$A \subseteq B$且$B \subseteq A$则$A = B$
\end{itemize}
\subsubsection{真子集}
\begin{itemize}
\item 反自反性：对任意集合$A$,$A \subsetneq A$为假。
\item 传递性：若$A \subsetneq B$且$B \subsetneq C$,则$A \subsetneq C$
\item 反对称性：若$A \subsetneq B$,则$B \subsetneq A$为假。
\end{itemize}
\subsection{符号 ⊂ 和 ⊃}
有些作者使用符号$\subset$和$\supset$分别表示子集与超集；也就是说，它们的含义与符号$\subseteq $和$\supseteq$相同并可以互换。例如，在这些作者看来，对于任意集合$A$，都有$A \subset A$，这是一个自反关系。

另一些作者则倾向于使用符号$\subset$和$\quad \supset$分别表示真子集（也称严格子集）与真超集；也就是说，它们的含义与符号$\subsetneq $和$\supsetneq$相同并可以互换。这种用法使得$\subseteq$与$\subset$类似于不等号$\leq $与$ <$例如，如果$x \leq y$，那么$x$可能等于$y$，也可能不等于；但如果$x < y$，那么$x$一定不等于$y$，而且小于$y$（这是一个反自反关系）。同样地，在采用“⊂ 表示真子集”的约定下：如果$A \subseteq B$，那么$A$可能等于$B$，也可能不等于；但如果$A \subset B$，那么$A$一定不等于$B$。
\subsection{子集的例子}
\begin{figure}[ht]
\centering
\includegraphics[width=6cm]{./figures/bed3b9a86ef7ec5a.png}
\caption{正多边形构成多边形的一个子集。} \label{fig_ZiJ_3}
\end{figure}
\begin{itemize}
\item 集合$A = {1, 2}$是集合$B = {1, 2, 3}$的真子集，因此表达式
$A \subseteq B(A \subseteq B)$和$A \subsetneq B(A \subsetneq B)$都成立。
\item 集合$D = {1, 2, 3}$是集合$E = {1, 2, 3}$的子集（但不是真子集），因此$D \subseteq E(D \subseteq E)$成立，而$D \subsetneq E(D \subsetneq E)$不成立（即为假）。
\item 集合{x : x 是大于 10 的素数}是集合{x : x 是大于 10 的奇数}的真子集。
\item 自然数集合是有理数集合的真子集；同样，线段上的点集是直线上点集的真子集。这是两个例子，其中子集与整体集合都是无限的，并且子集与整体集合具有相同的基数（对应于有限集的大小，即元素的数量）；这种情况可能会违背人们最初的直觉。
\item 有理数集合是真实数集合的真子集。在这个例子中，这两个集合都是无限的，但后者集合的基数（或势）大于前者集合。
\end{itemize}

另一个在欧拉图中的例子：
\begin{figure}[ht]
\centering
\includegraphics[width=6cm]{./figures/f9de8dcdb02e9233.png}
\caption{A 是 B 的真子集。} \label{fig_ZiJ_4}
\end{figure}
\begin{figure}[ht]
\centering
\includegraphics[width=6cm]{./figures/9d616e38414a8b0c.png}
\caption{翻译如下：C 是 B 的子集，但不是 B 的真子集。} \label{fig_ZiJ_5}
\end{figure}
\subsection{幂集}
集合$S$的所有子集所组成的集合称为它的幂集，记作$\mathcal{P}(S)$.

包含关系 $\subseteq$ 在幂集 $\mathcal{P}(S)$ 上定义了一个偏序：$A \leq B \iff A \subseteq B$.

我们也可以用“反向包含”来在$\mathcal{P}(S)$ 上定义偏序：$A \leq B $当且仅当$B \subseteq A$.对于集合 $S$ 的幂集 $\mathcal{P}(S)$，其包含关系的偏序同构于$|S|=k$个$\{0,1\}$上偏序的笛卡尔积，其中$\{0,1\}$的偏序为 $0<1$。这可以通过如下方式说明：
将集合

$$
S = \{s_{1}, s_{2}, \ldots, s_{k}\}
$$

枚举，并把每个子集 $T \subseteq S$（即 $T \in 2^{S}$）对应到一个 $\{0,1\}^{k}$ 的 k 元组：其中第 i 个坐标为 1 当且仅当 $s_i \in T$。

集合 $A$ 的所有 k 元子集记作

$$
\binom{A}{k},
$$

这一记号与二项式系数的记法类似，因为二项式系数表示从一个 n 元集合中选取 k 元子集的个数。在集合论中，记号

$$
[A]^k
$$

也很常见，特别是在 $k$ 是超限基数时。

---

### 包含关系的其他性质

* **交集刻画**：
  集合 $A$ 是集合 $B$ 的子集，当且仅当它们的交集等于 $A$。形式化地：

$$
A \subseteq B \iff A \cap B = A.
$$

* **并集刻画**：
  集合 $A$ 是集合 $B$ 的子集，当且仅当它们的并集等于 $B$。形式化地：

$$
A \subseteq B \iff A \cup B = B.
$$

* **有限集的基数刻画**：
  若 $A$ 是有限集，则 $A \subseteq B$ 当且仅当

$$
|A \cap B| = |A|.
$$

---

### 子集关系与布尔代数

子集关系在集合上定义了一个偏序。事实上，给定集合的所有子集在子集关系下构成一个布尔代数，其中：

* **并 (join)** 运算是交集，
* **交 (meet)** 运算是并集，
* 子集关系本身就是布尔代数中的包含关系。

---

### 包含是典范偏序

包含关系是典范的偏序，意思是：任意偏序集 $(X, \preceq)$ 都同构于某个按包含关系排列的集合族。

序数是一个简单的例子：如果把每个序数 $n$ 与集合

$$
[n] = \{ \text{所有小于等于 } n \text{ 的序数} \}
$$

对应起来，则有

$$
a \leq b \iff [a] \subseteq [b].
$$

---

要不要我帮你把这一段翻译整理成一个更像教材里的结构化讲解（加上小标题和公式展示），这样读起来更清晰？
